\section{工程实现与代码结构}

这一章我们来讲项目的具体工程实现，包括代码怎么组织的、核心模块怎么分工的，以及训练和评估的完整流程。

\subsection{项目目录结构}

为了让代码便于维护和扩展，我们把整个项目按功能拆成了几个独立的目录，具体结构如下：

\begin{lstlisting}[language=bash,caption={项目目录结构}]
project_root/
  src/
    data/          # 数据加载、预处理、CLIP特征缓存
    models/        # 模型定义(CLIP, A2 encoder, fusion, caption model)
    train/         # 训练脚本(baseline, fusion, alignment等)
    eval/          # 评估(generation, retrieval, sanity check, metrics)
    losses/        # 损失函数(contrastive, EEG augmentation, similarity)
    utils/         # 工具(config loader, seed, checkpoint, logger)
  scripts/         # 启动脚本和数据工具
  configs/         # YAML配置文件
  outputs/         # 训练输出(checkpoint, metrics, report)
  tests/           # 单元测试
\end{lstlisting}

\subsection{核心模块说明}

下面我们用表~\ref{tab:code-modules} 把几个重要模块的功能梳理一下。

\begin{table}[htbp]
  \centering
  \caption{核心代码模块}
  \label{tab:code-modules}
  \begin{tabularx}{\linewidth}{l X}
    \toprule
    模块 & 功能 \\
    \midrule
    \texttt{src/data/dataset.py} & EEGVisionCaptionDataset，支持 JSONL manifest 的数据加载器，兼容 Thought2Text、EEG-ImageNet、THINGS-EEG2 等格式 \\
    \texttt{src/data/clip\_cache.py} & CLIP 视觉特征预缓存模块 \\
    \texttt{src/models/eeg\_encoder.py} & 13 种 EEG encoder variant 定义及 \texttt{build\_eeg\_encoder} 工厂函数 \\
    \texttt{src/models/fusion.py} & GatedFusion 门控残差融合模块 \\
    \texttt{src/models/caption\_model.py} & SoftPromptCaptionModel 与内置 byte-level 回退 \\
    \texttt{src/train/train\_semantic\_fusion.py} & A2 语义融合训练脚本 \\
    \texttt{src/train/train\_align.py} & EEG-to-CLIP 对齐训练（InfoNCE+MSE）\\
    \texttt{src/eval/generate.py} & 多模式生成式评估 \\
    \texttt{src/eval/semantic\_fusion\_classifier\_eval.py} & A2 语义融合分类评估 \\
    \bottomrule
  \end{tabularx}
\end{table}

\subsection{训练与评估流程}

\textbf{训练流程}大致分六个步骤：(1) 先做数据准备——构建 JSONL manifest，然后生成 CLIP 特征缓存；(2) 可选阶段，做 A2 EEG encoder 的对齐训练——用 InfoNCE 对比损失把 EEG 特征映射到 CLIP 空间；(3) A2 语义融合训练——把 fusion 和 classification 模块放在一起联合训练；(4) VTF token 融合训练；(5) 生成式模型训练——这里用的是 Qwen2-VL + LoRA；(6) 最后做全量测试——把所有退化条件 $\times$ EEG 模式的组合都跑一遍。

\textbf{输出文件}方面，每次训练的结果都统一放在 \texttt{outputs/<run\_name>/} 下面，包括最佳 checkpoint、训练日志、metrics JSON 和 report。生成式评估会输出 JSONL 格式的文件，每条记录里有 image\_id、mode、reference、prediction 这几个字段。

\subsection{遇到的工程问题}

开发过程中我们碰到了一些实际问题，也逐一解决了：(1) 增强视觉 token 的维度跟 Qwen2-VL 的输入要求不完全对得上，所以我们自己在 VTF 后面加了一个轻量的 prefix adapter，专门做格式对齐和维度映射；(2) 部分 BLIP 生成的伪 caption 有重复、太长或者跟类别对不上的情况，我们通过最小长度过滤（$\geq$ 3 词）加正则去噪来处理；(3) 全量测试要跑 1997 $\times$ 6 $\times$ 4 $\approx$ 48000 次推理，生成式模型的自回归解码是最耗时的环节，完整跑一轮大概需要 2--4 小时；(4) 主要的训练任务都是在单张 48 GB GPU 上完成的，没有用 DeepSpeed 或者多 GPU 分布式训练。
